\SourceLeaf{298}
Aus $\dfrac{6^\ast\times 8^\ast}{11^\ast}$ oder $\dfrac{6^\ast\times 10^\ast}{9^\ast}$ oder $\dfrac{8^\ast\times 10^\ast}{7^\ast}$ folgt

\[
PP=qq,
\]

woraus notwendig

\[
P=+q
\]

geschlossen werden muss, sofern man nur Dreiecke betrachtet, deren Seiten und Winkel nicht über $180^0$ hinausgehen. Dann folgt unmittelbar

\[
Q=+s,\qquad R=+p,\qquad S=+r,
\]

und dies sind zusammen die vier gewünschten Sätze.

Ich überlasse es Ihnen, inwieweit Sie hiervon oder von anderen mir geläufigen Anordnungen der Rechnung in Ihrem Werk Gebrauch machen möchten. \mbox{.\,.\,.}

\medskip

\Needspace{8\baselineskip}
\begin{center}
[2.]

\textsc{Gauss} an \textsc{Gerling}. Göttingen, 26. Juni 1816.
\end{center}

\mbox{.\,.\,.\,.\,.\,.\,.} Ich war bisher durch so viele andere Dinge abgelenkt, dass ich noch nicht über die zweckmäßigste Reihenfolge der ersten Sätze der sphärischen Trigonometrie bei einer geometrischen Behandlung nachdenken konnte. Ich meine jedoch, am einfachsten wäre es, Euklid XI. 20. 21 zugrunde zu legen und den Satz, dass $A+B+C>180^0$ ist, aus dem Polardreieck abzuleiten. Ich glaube kaum, dass die Betrachtung des ebenen Dreiecks zwischen den Eckpunkten des sphärischen Dreiecks einen erheblichen Gewinn bringt, da es zwischen dessen Winkeln keine \textls{sehr} einfachen Beziehungen zu geben scheint. Die folgende ist wohl eine der elegantesten. Sind $A'$, $B'$, $C'$ die Winkel des ebenen Dreiecks und gilt

\[
A+B+C-180^0=u,
\]

so ist

\[
\cos A'=\cos\tfrac{1}{2}a\cos(A-\tfrac{1}{2}u).
\]

\mbox{.\,.\,.\,.\,.\,.\,.\,.\,.}
